\documentclass[../../../include/open-logic-section]{subfiles}

\begin{document}

\olfileid{sth}{story}{extensionality}
\olsection{సమితుల సమానత్వ సూత్రం (Extensionality)}

మొదట చెప్పాల్సిన విషయం:
సమితుల గుర్తింపు వాటి
\tetoken{మూలకాల}{!!{element}s}
ద్వారానే నిర్ణయించబడుతుంది.
మరింత ఖచ్చితంగా:

\begin{axiom}[సమితుల సమానత్వ సూత్రం]
$A$, $B$ సమితులకు ఒకే
\tetoken{మూలకాలు}{!!{element}s}
ఉంటే, $A$, $B$ ఒకే సమితి.
\[
  \lforall[A][\lforall[B][(\lforall[x][(x \in A \liff x \in B)] \lif
  \eq[A][B])]]
\]
\end{axiom}

\olref[sfr][][]{part} అంతటా దీనిని
పరికల్పించాం. అయినా మళ్లీ చెప్పడం
ముఖ్యం. సమితుల సమానత్వ స్వీకృతం,
సమితి దాని \tetoken{మూలకాల}{!!{element}s}
ద్వారా నిర్ణయించబడుతుందనే
మూలభావాన్ని వ్యక్తపరుస్తుంది.
(ఇక్కడ సమితులను \emph{భావనలతో}
పోల్చవచ్చు: సరిగ్గా అవే వస్తువులు
అనేక భిన్న భావనల కిందకు రావచ్చు.)

ఈ సూత్రాన్ని ఎందుకు స్వీకరించాలి?
దీనిని నిరాకరించడం అంటే
\emph{సమితి సిద్ధాంతం} అని
పిలవదగిన దేనినైనా వదిలేయడమే
అని చెప్పడం సమంజసం.
సమితి సిద్ధాంతం అనేది
మూలకాధారిత సముదాయాల
సిద్ధాంతమే; అంతకంటే
ఎక్కువా తక్కువా కాదు.

అయితే \olref[sth][][]{part}లో
నిజమైన సవాలు \emph{ఏ సమితులు
ఉన్నాయి} అని చెప్పే సూత్రాలను
నిర్దేశించడం. ఈ ప్రశ్నకు
నిజంగా ``స్పష్టమైనది'' అనిపించే
ఒకే జవాబు నిరూపితంగా తప్పని
తేలుతుంది.

\end{document}
